% 71-03(71쪽) — 보기 ①~⑤ 의 그래프 다섯 개(닫힌구간 [0,2] 에서의 y=f(x)). 보기 자체가 그림이라 본문 아래에 둔다.
% 스캔 화질이 낮아 TikZ 로 다시 그림 — 라벨은 원문 그대로(O, 2, x, y, y=f(x))만 두고 원문에 없는 값은 넣지 않는다.
\begin{tikzpicture}[x=0.4cm, y=0.4cm, line width=0.5pt, font=\footnotesize]
  % ① 상수함수
  \begin{scope}
    \node at (-1.9,2.2) {①};
    \draw[->, >=stealth, line width=0.7pt] (-0.9,0) -- (2.9,0) node[below=1pt] {$x$};
    \draw[->, >=stealth, line width=0.7pt] (0,-1.0) -- (0,2.7) node[left=1pt] {$y$};
    \node[below left=-2pt and -3pt] at (0,0) {$\pt{O}$};
    \draw[line width=0.6pt] (0,1.5) -- (2,1.5);
    \draw[dotted, line width=0.4pt] (2,1.5) -- (2,0);
    \node[below=1pt] at (2,0) {$2$};
    \node[anchor=west] at (0.8,2.15) {$y=f(x)$};
  \end{scope}
  % ② 아래로 볼록한 포물선 (x축 아래로 내려감)
  \begin{scope}[shift={(8.4,0)}]
    \node at (-1.9,2.2) {②};
    \draw[->, >=stealth, line width=0.7pt] (-0.9,0) -- (2.9,0) node[below=1pt] {$x$};
    \draw[->, >=stealth, line width=0.7pt] (0,-1.0) -- (0,2.7) node[left=1pt] {$y$};
    \node[below left=-2pt and -3pt] at (0,0) {$\pt{O}$};
    \draw[densely dashed, line width=0.4pt] (0,1.2) -- (2,1.2);
    \draw[dotted, line width=0.4pt] (2,1.2) -- (2,0);
    \draw[line width=0.6pt, domain=0:2, samples=80, smooth] plot (\x, {1.6*(\x-1)*(\x-1)-0.4});
    \node[below=1pt] at (2,0) {$2$};
    \node[anchor=west] at (0.8,2.15) {$y=f(x)$};
  \end{scope}
  % ③ V 꼴(꺾인 점에서 미분가능하지 않다)
  \begin{scope}[shift={(16.8,0)}]
    \node at (-1.9,2.2) {③};
    \draw[->, >=stealth, line width=0.7pt] (-0.9,0) -- (2.9,0) node[below=1pt] {$x$};
    \draw[->, >=stealth, line width=0.7pt] (0,-1.0) -- (0,2.7) node[left=1pt] {$y$};
    \node[below left=-2pt and -3pt] at (0,0) {$\pt{O}$};
    \draw[densely dashed, line width=0.4pt] (0,1.2) -- (2,1.2);
    \draw[dotted, line width=0.4pt] (2,1.2) -- (2,0);
    \draw[line width=0.6pt] (0,1.2) -- (1.15,0) -- (2,1.2);
    \node[below=1pt] at (2,0) {$2$};
    \node[anchor=west] at (0.8,2.15) {$y=f(x)$};
  \end{scope}
  % ④ 극대 뒤 극소가 x축 아래에 있는 곡선
  \begin{scope}[shift={(0,-4.6)}]
    \node at (-1.9,2.2) {④};
    \draw[->, >=stealth, line width=0.7pt] (-0.9,0) -- (2.9,0) node[below=1pt] {$x$};
    \draw[->, >=stealth, line width=0.7pt] (0,-1.0) -- (0,2.7) node[left=1pt] {$y$};
    \node[below left=-2pt and -3pt] at (0,0) {$\pt{O}$};
    \draw[densely dashed, line width=0.4pt] (0,0.7) -- (2,0.7);
    \draw[dotted, line width=0.4pt] (2,0.7) -- (2,0);
    \draw[line width=0.6pt, domain=0:2, samples=90, smooth] plot (\x, {0.7+1.1*sin(180*\x)});
    \node[below=1pt] at (2,0) {$2$};
    \node[anchor=west] at (1.0,1.6) {$y=f(x)$};
  \end{scope}
  % ⑤ 극소 뒤 극대가 모두 x축 위에 있는 곡선
  \begin{scope}[shift={(8.4,-4.6)}]
    \node at (-1.9,2.2) {⑤};
    \draw[->, >=stealth, line width=0.7pt] (-0.9,0) -- (2.9,0) node[below=1pt] {$x$};
    \draw[->, >=stealth, line width=0.7pt] (0,-1.0) -- (0,2.7) node[left=1pt] {$y$};
    \node[below left=-2pt and -3pt] at (0,0) {$\pt{O}$};
    \draw[densely dashed, line width=0.4pt] (0,0.9) -- (2,0.9);
    \draw[dotted, line width=0.4pt] (2,0.9) -- (2,0);
    \draw[line width=0.6pt, domain=0:2, samples=90, smooth] plot (\x, {0.9-0.55*sin(180*\x)});
    \node[below=1pt] at (2,0) {$2$};
    \node[anchor=west] at (0.8,2.15) {$y=f(x)$};
  \end{scope}
\end{tikzpicture}
